\documentclass[a4paper,12pt,subf,href,coursework]{disser8}
%%% master -- for master thesis
%%% bachelor -- for bachelor thesis
%%% specialist -- for specialist diploma
%%% coursework -- for course work

\usepackage{multirow} %объединение рядов по вертикале в таблицах
%\usepackage{hyperref}
\usepackage{comment}	
\usepackage{mathtools}
\usepackage{cmap} %делает возможным поиск и выделение символов в pdf
\usepackage{longtable} %длинные таблицы
\usepackage{array} % матрицы
%\usepackage{subfigure}
\DeclareUnicodeCharacter{2061}{}
\usepackage{amsmath}
\usepackage{cases}


\usepackage{paralist} %пакет для невложенных нумерованных и ненумерованных списков, оформленных по требованиям госта
\usepackage{enumitem} %если списки вложенные, то нужен вот этот пакет и три следующие за ним команды
\AddEnumerateCounter{\asbuk}{\@asbuk}{щ}
\setlist[1]{leftmargin=0cm,itemindent=\parindent}
\setlist[2]{leftmargin=\parindent,listparindent=\parindent,itemindent=\parindent}

\usepackage{subcaption, floatrow, graphicx, calc} % картинки
\floatsetup{floatrowsep=qquad}
\usepackage{epstopdf} % пакет для автоконвертации eps и ps картинок в pdf
\epstopdfsetup{update,suffix=,prefersuffix=false}
\AppendGraphicsExtensions{.ps}

\usepackage[paperwidth=21cm,paperheight=29.7cm,left=3cm,right=2cm,top=2cm,bottom=2cm]{geometry} % определение макета страницы
\usepackage[T2A]{fontenc}
\usepackage[utf8]{inputenc} % кодировка utf-8
\usepackage[english,russian]{babel} % переносы для русского и английского языков
%\usepackage{subfig} %пакет для картинок в картинке
\usepackage{wasysym} %всякие символы
\usepackage{caption} %подписи к таблицам и рисункам
\captionsetup[figure]{labelsep=endash,justification=justified,singlelinecheck=false,font=small}
\captionsetup[table]{labelsep=endash,justification=justified,singlelinecheck=false,font=normalsize}

% правильная нумерация источников в списке литературы
\makeatletter
\renewcommand{\@biblabel}[1]{#1}
\makeatother

%нумерация разделов приложения заглавными русскими буквами
\renewcommand\theappendix{\Asbuk{chapter}}
\setcounter{tocdepth}{3}
\renewcommand{\topfraction}{0.85}
\renewcommand{\textfraction}{0.1}
\renewcommand{\floatpagefraction}{0.75}
\renewcommand{\theenumi}{\asbuk{enumi}}
\renewcommand{\theenumii}{\arabic{enumii}}
\begin{document}

%%%%%% Общие поля титульного листа %%%%%%
\institution{\textbf{Министерство науки и высшего образования Российской Федерации\\
ФГАОУ ВО <<УрФУ имени первого Президента России Б. Н. Ельцина>>}}
%\department{наук о Земле} % Название департамента 
\chair{астрономии, геодезии, экологии и мониторинга окружающей среды} % Название кафедры
\topic{ТЕМА} % Тема
\city{Екатеринбург} % Год, дата

%%%%%% Титульный лист курсовой работы %%%%%%
\title{Тема задания на практику\\ТЕМА\\ОТЧЕТ\\ 
Вид практики Производственная практика
Тип практики Практика по получению профессиональных умений и опыта
профессиональной деятельности} % Тип
\groupnumber{МЕН-} % Номер группы
\author{АВТОР} % ФИО автора в именительном падеже
\gradcode{03.05.01 Астрономия}
\sa {ОРГ}{ПЕРВЫЙ РУКОВОДИТЕЛЬ}
\sasecond {УРФУ}{ВТОРОЙ РУКОВОДИТЕЛЬ}
 %В этом файле лежат все команды для титульного листа

\def\apj{Astrophys.~J}
\def\aatr{Astron.~Astroph.~Trans}
\def\aaps{Astron.~and Astrophys.~Suppl.~Ser}
\def\pasp{Publ.~Astron.~Soc.~Pac}
\def\gca{Geochim.~Cosmochim.~Acta}
\def\aap{Astron.~Astrophys}
\def\aspcs{ASP~Conf.~Ser}
\def\asrep{Astron.~Rep}
\def\nat{Nature}
\def\apjl{Astrophys.~J.~Lett}
\def\apjs{Astrophys.~J.~Suppl.~Ser}
\def\aj{Astron.~J}
\def\mnras{Mon.~Not.~R.~Astron.~Soc}
\def\araa{Ann.~Rev.~Astron.~Astrophys}
\def\jcp{J.~Chem.~Phys}
\def\apss{Astrophys.~Space.~Sci}
\def\prl{Phys.~Rev.~Lett}
\def\phrva{Phys.~Rev.~A}
\def\phlb{Phys.~Let.~B}
\def\pf{Phys.~Fluids}
\def\azh{Астрон.~журн}
\def\pazh{Письма~в~Астрон.~журн}
\def\jgr{J.~Geophys.~Res}
\def\cemda{Celest.~Mech.~Dyn.~Astr}
\def\jcoph{J.~Comp.~Phys}
\def\cophc{Comput.~Phys.~Commun}
\def\phpl{Physics~of~Plasmas}
\def\pasj{Publ.~Astron.~Soc.~Jpn}
\def\avest{Астрон.~вестн}
\def\jrasc{J.~R.~Astron.~Soc.~Can}
\def\cemec{Celest.~Mech}
\def\pasau{Proc.~Astron.~Soc.~Aust}
\def\puasau{Publ.~Astron.~Soc.~Aust}
\def\jasa{J.~Acoust.~Soc.~Am}
\def\jfm{J.~Fluid~Mech}
\def\cajph{Can.~J.~Phys}
\def\mitag{Mitt.~Astron.~Ges}
\def\bain{Bull.~Astron.~Inst.~Neth}
\def\epsl{Earth~Planet.~Sci.~Lett}
\def\ibvs{Inf.~Bull.~Variable~Stars}
\def\arep{Astr.~Rep}
\def\phr{Phys.~Rep}
\def\astl{Astron.~Letters}
\def\sci{Science}
\def\jqsrt{J.~Quant.~Spectrosc.~Radiat.~Transfer}
\def\emp{Earth,~Moon~and~Planets}
\def\icar{Icarus}
\def\pss{Planet.~Space~Sci}
\def\qjras{Q.~J.~R.~Astron.~Soc}
\def\nimpa{Nucl.~Instrum.~Methods~Phys.~Res.,~Sect.~A}
\def\soph{Sol.~Phys}
\def\lnm{Lect.~Notes~in~Math}
\def\an{Astron.~Nach}
\def\aph{Astroparticle~Physics}
\def\adspr{Adv.~Space~Res}
\def\geoj{Geophys.~J}
\def\caosp{Contrib.~Astron.~Obs.~Skalnat{\'e}~Pleso}
\def\vestcpbu{Вестн.~С.-Петерб.~ун-та}
\def\izvvrad{Изв.~вузов.~Радиофизика}
\def\izvans{Изв.~AH~CCCP}
\def\vestvgu{Вестн.~ВолГУ}
\def\vestsibgau{Вестн.~СибГАУ}
\def\bamass{Bull.~Am.~Astron.~Soc}
\def\rmxaa{Rev.~Mex.~Astron.~Astrofis}
\def\aapr{Astron.~Astrophys.~Rev}
\def\acp{Atmosphere~Chem.~Phys}
\def\cosiss{Космич.~исслед}
\def\ssrv{Space~Sci.~Rev}
%%%%%%%%%%%%
%%%%%%%%%%%%%%%%%%%%%%%
\def\jmph{J.~Math.~Phys}
\def\rvmps{Rev.~Mod.~Phys.~Suppl}
\def\rvmp{Rev.~Mod.~Phys}
\def\prd{Phys.~Rev.~D}
\def\nuphs{Nuc.~Phys.~B~Proc.~Suppl}
\def\nuphb{Nuc.~Phys.~B}
\def\skytel{Sky~Telesc}
%%%%%%%%%%%%
\def\thmc{Тез.~международ.~конф}
\def\mmc{Материалы~международ.~конф}
\def\mvrc{Материалы~всерос.~конф}
\def\cntc{Сб.~науч.~тр.~конф}
\def\tmnpc{Тр.~Международ.~науч.-практ.~конф}
\def\ctc{Сб.~тр.~конф}
\def\mcnctone{Тр.~31-й Международ.~студ.~науч.~конф., Екатеринбург, 28 янв.~--- 1~февр. 2002~г}
\def\mcnctto{Тр.~32-й Международ.~студ.~науч.~конф., Екатеринбург, 3---7 февр. 2003~г}
\def\mcncttr{Тр.~33-й Международ.~студ.~науч.~конф., Екатеринбург, 2---6 февр. 2004~г}
\def\mcnctfo{Тр.~34-й Международ.~студ.~науч.~конф., Екатеринбург, 31 янв.~--- 4~февр. 2005~г}
\def\mcnctfi{Тр.~35-й Международ.~студ.~науч.~конф., Екатеринбург, 30 янв.~--- 3~февр. 2006~г}
\def\mcnctsi{Тр.~36-й Международ.~студ.~науч.~конф., Екатеринбург, 29 янв.~--- 2~февр. 2007~г}
\def\mcnctse{Тр.~37-й Международ.~студ.~науч.~конф., Екатеринбург, 28 янв.~--- 1~февр. 2008~г}
\def\mcnctei{Тр.~38-й Международ.~студ.~науч.~конф., Екатеринбург, 2---6 февр. 2009~г}
\def\mcnctni{Тр.~39-й Международ.~студ.~науч.~конф., Екатеринбург, 1---5 февр. 2010~г}
\def\mcncforty{Тр.~40-й Международ.~студ.~науч.~конф., Екатеринбург, 31 янв.~--- 4~февр. 2011~г}
\def\mcncfone{Тр.~41-й Международ.~студ.~науч.~конф., Екатеринбург, 30 янв.~--- 3~февр. 2012~г}
\def\tmc{Тр.~Международ.~конф}
\def\tc{Тр.~конф}
\def\thc{Тез.~конф}

\def\IAUsympc{Proc.~IAU~Symp}
\def\IAUsymp{Proc.~IAU~Symp}
\def\IAUcoll{Proc.~IAU~Colloquia}
\def\prconf{Proc.~conf}
\def\princonf{Proc.~int.~conf}
\def\mvrnc{Материалы~всерос.~науч.~конф}
 %В этом файле команды для названий журналов
\maketitle

% ----------------------------------------------------------------
%\newpage
\tableofcontents
\include{designation} % раздел "ОБОЗНАЧЕНИЯ И СОКРАЩЕНИЯ"
% ----------------------------------------------------------------
\intro\label{ch:intro}
 % раздел "ВВЕДЕНИЕ"
\newpage
%\main % печатает заголовок "ОСНОВНАЯ ЧАСТЬ", не включается в курсовую работу

\chapter{Обзор литературы}\label{ch:review}
 % раздел "Обзор литературы"
%\newpage

\chapter{Методика}\label{ch:metods}
 % раздел "Способы и методы теоретических и аналитических исследований"
\newpage


\chapter{Результаты}
 % раздел "Результаты и их обсуждение"
\newpage


\conclusion\label{ch:concl}
 % раздел "ЗАКЛЮЧЕНИЕ"
\newpage


\bibliographystyle{Kourovkastyle} %Стилевой файл для списка литературы в BibTeX
\bibliography{ref} % подключение списка литературы, сделанного в формате .bib
\newpage
% ----------------------------------------------------------------
%\appendix % Раздел с Приложениями

%\chapter*{ПРИЛОЖЕНИЕ А}
\addcontentsline{toc}{chapter}{ПРИЛОЖЕНИЕ А}

\chapter*{ПРИЛОЖЕНИЕ Б}
\addcontentsline{toc}{chapter}{ПРИЛОЖЕНИЕ Б} % Файл с текстом приложения

\end{document}